\textbf{Задача 25} Приведите пример языка лежащего в $\bpp$, для которого неизвестно доказательство принадлежности ни $\rp$, ни $\corp$.

\textbf{Лемма Шварца-Зиппеля}
Пусть задан ненулевой многочлен $S$ от $n$ переменных, степени $\le d$ над $\mathds{F}_p$. Тогда
\[\mathrm{Pr}\left\{S(\vec x)\neq0\right\}\ge 1 - \frac{dn}{p}\]
\begin{proof} по индукции. База $n=1$ --- у многочлена степени $d$ будет $\le d$ корней. Значит $\ge p - d$ аргументов с не нулевым значением, а всего аргументов $p$.
Значит
\[\mathrm{Pr}\left\{S(\vec x)\neq0\right\}\ge\frac{p-d}{p}= 1 - \frac{d}{p}\]

Переход

Разложим $S$ по степеням $x_{n+1}$
\[S(\vec x) = P_0(\vec x) + x_{n + 1}P_2(\vec x) + x_{n + 1}^2 P_1(\vec x) + \ldots + x_{n+1}^d P_d(\vec x)\quad (P_j(\vec x)\not \equiv 0 \text{ для некоторого } j)\]

Сумма может обнуляться когда все коэффициенты нулевые, т.~е. $(x_1,\ldots,x_n)$ --- корень $P_0$, $P_1$, \ldots, $P_d$. То есть как минимум корень $P_j$, что по предположению индукции происходит с вероятностью $\le\frac{dn}{p}$.

Так же сумма может обнуляться когда $x_{n+1}$ корень многочлена, но по предыдущим рассуждениям это происходит с вероятностью $\le\frac{d}{p}$. Значит итоговая вероятность стать нулём $\le\frac{d(n+1)}{p}$ \end{proof}

\textbf{Решение задачи}
Определим $A=\{(P,Q,R)\ |\ P,Q,R$ --- многочлены в сжатой записи, и $2$ из них равны между собой и отличаются от третьего$\}$.
Покажем, что $A\in\bpp$

Возьмем достаточно большое просто число $p$ (зависит от входа), возьмем случайных $\vec x$ и проверим, что $P(\vec x)=Q(\vec x)$, результат равенства положим в переменную $\xi_{p,q}$. Аналогично найдем $\xi_{p,r}$ и $\xi_{q,r}$. Вернем $1$, если среди $(\xi_{p,q},\xi_{p,r},\xi_{q,r})$ ровно два нуля и одна единица. Иначе вернем ноль. 

Обозначим этот алгоритм за $M$. Тогда
\[\forall (P,Q,R)\in A\quad \mathrm{Pr}(M(P,Q,R) = 1)\ge 1 \cdot \left(1 - \frac{dn}{p}\right)^2 \ge 1 - \left(\frac{dn}{p}+\frac{dn}{p}\right) \ge \frac{2}{3}\]

Первое неравенство следует из того, что у нас в любом случае будет хотя бы одна единичка, но также могут появиться лишние.


\[\forall (P,Q,R)\not\in A\quad \mathrm{Pr}(M(P,Q,R) = 1)\le 3 \cdot \frac{dn}{p} \le \frac{1}{3}\]

Предпоследнее неравенство следует из того, что в случае если они все равны между собой, то мы всегда вернем ноль, а если среди них нет равных, то хотя бы одно из $(\xi_{p,q},\xi_{p,r},\xi_{q,r})$ должно вернуть $1$, что я и написал.

В итоге $A\in \bpp$. Про принадлежность к $\rp$ или к $\corp$ мне не известно.
